\hypertarget{python-le-langage}{%
\section{Python, le langage}\label{python-le-langage}}

Le langage de programmation que vous apprendrez est Python.

\begin{description}
\item[digraph CPython\_Architecture \{]
// Disposition horizontale rankdir=LR;

// Style par défaut des nœuds node {[}shape=box, style=filled,
fillcolor="\#f8f9fa", color="\#343a40", fontname="Arial"{]};

// Définition des nœuds Source {[}label="Code Sourcen(.py)",
fillcolor="\#e3f2fd"{]}; Parser {[}label="Analyseurn(lexer \&
parser)"{]}; AST {[}label="Arbre syntaxiquenabstrait (AST)"{]}; Compiler
{[}label="Compilateur"{]}; Bytecode {[}label="Bytecode"{]}; PVM
{[}label="Machine VirtuellenPython", fillcolor="\#fff3cd"{]}; Output
{[}label="RésultatnExécuté", shape=oval, fillcolor="\#d4edda"{]};

// Connexions Source -\textgreater{} Parser {[}label=" Texte"{]}; Parser
-\textgreater{} AST {[}label=" Tokens"{]}; AST -\textgreater{} Compiler;
Compiler -\textgreater{} Bytecode; Bytecode -\textgreater{} PVM; PVM
-\textgreater{} Output;
\end{description}

\}

Il y a deux façons d\textquotesingle utiliser le Interprète: mode
interactif et mode script. Dans mode interactif, vous tapez des
programmes Python et l\textquotesingle interpréteur affiche le résultat:

\begin{Shaded}
\begin{Highlighting}[]
\OperatorTok{\textgreater{}\textgreater{}\textgreater{}} \DecValTok{1} \OperatorTok{+} \DecValTok{1}
\DecValTok{2}
\end{Highlighting}
\end{Shaded}

Les chevrons, \texttt{\textgreater{}\textgreater{}\textgreater{}}, sont
l\textquotesingle{}\textbf{invite} de commande que
l\textquotesingle interprèteur utilise pour indiquer
qu\textquotesingle il est prêt. Si vous tapez \texttt{1\ +\ 1}, Python
répond 2.

Alternativement, vous pouvez stocker du code dans un fichier et utiliser
l\textquotesingle interpréteur pour exécuter le contenu du fichier, qui
est appelé un \textbf{script}. Par convention, les scripts Python ont
des noms qui se terminent par {.py}.

Pour exécuter le script, vous devez indiquer à
l\textquotesingle interpréteur l\textquotesingle emplacement du fichier.
Si vous avez un script nommé \texttt{monscript.py}, vous
l\textquotesingle exécutez avec la commande
\texttt{python\ monscript.py}.

Travailler en mode interactif est pratique pour tester de petits
morceaux de code car vous pouvez les taper et les exécuter
immédiatement. Pour quoi que ce soit de plus que quelques lignes, vous
devriez enregistrer votre code en tant que script pour que vous puissiez
le modifier et l\textquotesingle exécuter plus tard.

\hypertarget{langages-naturels-et-formels}{%
\subsection{Langages naturels et
formels}\label{langages-naturels-et-formels}}

Les langues naturelles sont les langues que les gens parlent, comme
l\textquotesingle anglais, l\textquotesingle espagnol et le français.
Ils n\textquotesingle ont pas été conçus par les gens (bien que les gens
essaient de leur imposer un ordre); Ils ont évolué naturellement.

Les langues formelles sont des langues conçues par des personnes pour
applications spécifiques. Par exemple, la notation que les
mathématiciens L\textquotesingle utilisation est un langage formel qui
est particulièrement bon pour désigner relations entre les nombres et
les symboles. Les chimistes utilisent un formel langage pour représenter
la structure chimique des molécules. Et Le plus important:
